% Raport z Etapu 3 — Flashgrad (stan gałęzi main, commit 633007b)
% Kompilacja: pdflatex raport.tex (dwa przebiegi) lub Overleaf (pdfLaTeX).
\documentclass[11pt,a4paper]{article}

\usepackage[utf8]{inputenc}
\usepackage[T1]{fontenc}
\usepackage[polish]{babel}
\usepackage{lmodern}
\usepackage[margin=2.3cm]{geometry}
\usepackage{microtype}
\usepackage{booktabs}
\usepackage{tabularx}
\usepackage{array}
\usepackage{xcolor}
\usepackage{float}
\usepackage{caption}
\usepackage{enumitem}
\usepackage{listings}
\usepackage{amsmath}
\usepackage{siunitx}
\usepackage{tikz}
\usetikzlibrary{arrows.meta,positioning,shapes.geometric,calc,babel}
\usepackage{pgfplots}
\pgfplotsset{compat=1.18}
\usepackage{fancyhdr}
\usepackage[hidelinks]{hyperref}

\sisetup{output-decimal-marker={,},group-separator={\,}}

\definecolor{fgblue}{HTML}{2A6FDB}
\definecolor{fggreen}{HTML}{2E9E5B}
\definecolor{fgorange}{HTML}{E08A1E}
\definecolor{fgred}{HTML}{C8423B}
\definecolor{fgpurple}{HTML}{7C4DCC}
\definecolor{fggray}{HTML}{6B7280}
\definecolor{codebg}{HTML}{F6F8FA}

\lstdefinestyle{term}{
  basicstyle=\ttfamily\footnotesize,
  backgroundcolor=\color{codebg},
  frame=single, rulecolor=\color{fggray!40},
  columns=fullflexible, keepspaces=true,
  breaklines=true, showstringspaces=false,
  inputencoding=utf8, extendedchars=true,
  literate={ą}{{\k{a}}}1 {ć}{{\'c}}1 {ę}{{\k{e}}}1 {ł}{{\l{}}}1 {ń}{{\'n}}1
           {ó}{{\'o}}1 {ś}{{\'s}}1 {ź}{{\'z}}1 {ż}{{\.z}}1
}
\lstdefinestyle{cpp}{
  style=term, language=C++,
  keywordstyle=\color{fgblue}\bfseries,
  commentstyle=\color{fggray}\itshape,
  stringstyle=\color{fggreen}
}

\setlength{\headheight}{14pt}
\pagestyle{fancy}
\fancyhf{}
\lhead{\small Flashgrad -- Etap 3}
\rhead{\small Wydział MiNI, PW}
\cfoot{\thepage}

\tikzset{
  box/.style={draw, rounded corners=3pt, align=center, font=\small, minimum height=8mm,
              inner sep=4pt},
  layer/.style={box, minimum width=14cm, fill=#1!12, draw=#1!70!black},
  mod/.style={box, fill=#1!15, draw=#1!70!black, minimum width=2.9cm},
  arr/.style={-{Stealth[length=2.2mm]}, thick},
  darr/.style={-{Stealth[length=2.2mm]}, thick, dashed},
}

\title{\textbf{Dokument Etapu 3:\\ Implementacja rdzenia systemu i raport postępu}\\[4pt]
  \large Projekt inżynierski -- Wydział Matematyki i Nauk Informacyjnych, Politechnika Warszawska\\[2pt]
  \normalsize \textit{Flashgrad -- biblioteka uczenia głębokiego w C++20 i CUDA}}
\author{Konrad Burdach \and Wojciech Królikowski \and Antoni Rakowski}
\date{18 listopada 2026}

\begin{document}
\maketitle
\thispagestyle{fancy}

\begin{abstract}
\noindent
Dokument podsumowuje stan biblioteki \textbf{Flashgrad} na gałęzi \texttt{main}
(commit \texttt{633007b}) na koniec Etapu~3. Flashgrad to napisana od podstaw biblioteka
uczenia głębokiego w~C++20 i~CUDA, wzorowana na PyTorch: tensory jako widoki z~krokami,
własne jądra GPU, buforujący alokator pamięci, automatyczne różniczkowanie, moduły
\texttt{nn}, optymalizatory i~wielowątkowy \texttt{DataLoader}. Zrealizowano fazy 0--6
harmonogramu. Kamieniem milowym jest pełny przepływ end-to-end: uczenie perceptronu
wielowarstwowego na zbiorze MNIST wyłącznie na własnych jądrach, z~dokładnością testową
97,5\% (PyTorch z~tą samą konfiguracją: 97,7\%). Wszystkie 974 uruchomienia testów
(487 testów w~dwóch konfiguracjach alokatora) przechodzą, pokrycie kodu hosta wynosi
96,0\%, a~\texttt{compute-sanitizer} nie zgłasza błędów pamięci, wyścigów ani błędów
synchronizacji.
\end{abstract}

\tableofcontents
\newpage

%=====================================================================
\section{Postęp implementacji i stan repozytorium}
%=====================================================================

\begin{table}[H]
\centering
\caption{Stan repozytorium (gałąź \texttt{main}, 18.11.2026).}
\label{tab:repo}
\begin{tabularx}{\textwidth}{@{}lX@{}}
\toprule
\textbf{Wskaźnik} & \textbf{Wartość} \\
\midrule
Gałąź \texttt{main} & Commit \texttt{633007b} (PR \#17, 16.11.2026). Stabilna: pełny zestaw testów
  przechodzi w~kompilacji Release i~Debug, z~alokatorem buforującym i~bez niego \\
Liczba commitów & 92 (75 merytorycznych + 17 scalających) \\
Scalone Pull Requesty & 15 (\#1--\#4, \#6, \#8--\#17) \\
Udział w~commitach & K.~Burdach 64, W.~Królikowski 20, A.~Rakowski 8 \\
Rozmiar kodu & 5032 linie biblioteki (nagłówki + \texttt{src/}), 5634 linie testów,
  382 linie przykładów, 287 linii benchmarków; 133 pliki \\
Testy jednostkowe & 487 testów w~64 zestawach (GoogleTest), w~CTest 974 uruchomienia \\
Pokrycie kodu (host) & 96,0\% linii w~\texttt{src/} (gcov) \\
Środowisko & RTX 3090 (\texttt{sm\_86}, 24~GiB), CUDA 13.0, GCC 13.3, CMake 3.28, Ubuntu 24.04 \\
\bottomrule
\end{tabularx}
\end{table}

Projekt nie ma jeszcze zdalnego CI -- darmowe maszyny GitHub Actions nie mają GPU. Testy
uruchamiane są lokalnie na docelowej karcie przed scaleniem każdego Pull Requesta.
Kompilacja kodu hosta używa \texttt{-Werror}, więc każde ostrzeżenie blokuje budowanie.

\begin{figure}[H]
\centering
\begin{tikzpicture}
\begin{axis}[
  ybar, width=15cm, height=5.2cm, bar width=11pt,
  ymin=0, ymax=48,
  symbolic x coords={13.10,15.10,16.10,17.10,18.10,19.10,20.10,21.10,24.10,29.10,30.10,
                     31.10,11.11,13.11,14.11,15.11,16.11},
  xtick=data, ylabel={liczba commitów},
  nodes near coords, nodes near coords style={font=\scriptsize},
  x tick label style={font=\scriptsize, rotate=45, anchor=east},
  axis lines*=left, ymajorgrids, grid style={gray!20}]
\addplot[fill=fgblue!70, draw=fgblue!80!black] coordinates {
  (13.10,4) (15.10,4) (16.10,4) (17.10,4) (18.10,1) (19.10,6) (20.10,41) (21.10,2)
  (24.10,1) (29.10,3) (30.10,1) (31.10,2) (11.11,1) (13.11,4) (14.11,3) (15.11,1)
  (16.11,10)};
\end{axis}
\end{tikzpicture}
\caption{Liczba commitów w~poszczególnych dniach (łącznie 92, z~commitami scalającymi).
  Skok 20.10 to autograd: każda różniczkowalna operacja ma osobny commit z~testem gradientu.}
\label{fig:commits}
\end{figure}

%=====================================================================
\section{Zrealizowane komponenty rdzenia}
%=====================================================================

\subsection{Architektura biblioteki}

Biblioteka ma budowę warstwową (rys.~\ref{fig:arch}). Użytkownik korzysta wyłącznie
z~API w~C++. Operacje matematyczne wykonują \emph{tylko} jądra CUDA -- tensory w~pamięci
CPU służą do tworzenia danych, wypisywania i~kopiowania. Pamięć GPU pochodzi z~własnego
alokatora buforującego, nigdy bezpośrednio z~\texttt{cudaMalloc}.

\begin{figure}[H]
\centering
\begin{tikzpicture}[node distance=3mm]
\node[layer=fgpurple, minimum height=15mm] (hi) {};
\node[anchor=north west, font=\small\bfseries] at ($(hi.north west)+(1mm,-1mm)$)
  {Warstwa wysoka (Fazy 5--6)};
\node[mod=fgpurple, anchor=south west] at ($(hi.south west)+(2mm,1.5mm)$) (nn)
  {\texttt{nn}: \texttt{Linear}, \texttt{Sequential},\\aktywacje, \texttt{Dropout}, straty};
\node[mod=fgpurple, right=3mm of nn, minimum width=3.6cm] (opt)
  {\texttt{optim}: \texttt{SGD},\\\texttt{Adam}/\texttt{AdamW}};
\node[mod=fgpurple, right=3mm of opt, minimum width=3.6cm] (data)
  {\texttt{data}: \texttt{Dataset},\\\texttt{DataLoader} (wątki)};
\node[layer=fggreen, below=of hi] (ag) {\textbf{Autograd} (\texttt{src/autograd})\quad
  węzły \texttt{Node}/\texttt{Edge}, \texttt{record()}, silnik \texttt{run\_backward}};
\node[layer=fgorange, below=of ag, minimum height=16mm] (core) {};
\node[anchor=north west, font=\small\bfseries] at ($(core.north west)+(1mm,-1mm)$)
  {Rdzeń (\texttt{src/core})};
\node[mod=fgorange, anchor=south west] at ($(core.south west)+(2mm,1.5mm)$) (shape)
  {\texttt{Shape}, \texttt{DType},\\ \texttt{Device}};
\node[mod=fgorange, right=2mm of shape] (stor) {\texttt{Storage}\\(licznik referencji)};
\node[mod=fgorange, right=2mm of stor] (tens) {\texttt{Tensor}\\(widok z~krokami)};
\node[mod=fgorange, right=2mm of tens] (ops) {dyspozytor ops\\+ walidacja};
\node[layer=fgred, below=of core, minimum height=16mm] (cuda) {};
\node[anchor=north west, font=\small\bfseries] at ($(cuda.north west)+(1mm,-1mm)$)
  {Backend CUDA (\texttt{src/cuda})};
\node[mod=fgred, anchor=south west] at ($(cuda.south west)+(2mm,1.5mm)$) (ew)
  {elementwise\\+ broadcast};
\node[mod=fgred, right=2mm of ew] (red) {redukcje,\\\texttt{log\_softmax}};
\node[mod=fgred, right=2mm of red] (mm) {\texttt{matmul}\\(register tiling)};
\node[mod=fgred, right=2mm of mm] (alloc) {alokator\\buforujący};
\node[layer=fggray, below=of cuda] (hw) {\textbf{CUDA Runtime 13.0} \quad
  $\rightarrow$ \quad GPU RTX 3090 (\texttt{sm\_86})};
\draw[arr] (hi) -- (ag);
\draw[arr] (ag) -- (core);
\draw[arr] (core) -- (cuda);
\draw[arr] (cuda) -- (hw);
\end{tikzpicture}
\caption{Warstwowa architektura biblioteki Flashgrad.}
\label{fig:arch}
\end{figure}

\subsection{Rdzeń tensorów (Faza 1)}

\texttt{Storage} to bufor pamięci ze zliczaniem referencji. \texttt{Tensor} to widok na
niego opisany trójką (\emph{kształt}, \emph{kroki}, \emph{przesunięcie}). Operacje
\texttt{t()}, \texttt{permute}, \texttt{slice}, \texttt{select}, \texttt{expand},
\texttt{view}, \texttt{squeeze}/\texttt{unsqueeze} tworzą nowy widok bez kopiowania
danych (rys.~\ref{fig:strides}); \texttt{reshape} kopiuje tylko wtedy, gdy nowego kształtu
nie da się wyrazić krokami. Dostępne są fabryki (\texttt{zeros}, \texttt{ones},
\texttt{arange}, \texttt{randn}), deterministyczny generator (\texttt{manual\_seed})
oraz kopie host $\leftrightarrow$ urządzenie.

\begin{figure}[H]
\centering
\begin{tikzpicture}[cell/.style={draw, minimum size=8mm, font=\small}]
\node[font=\small\bfseries, anchor=east] at (-0.5,0) {Storage:};
\foreach \i in {0,...,5} {
  \node[cell, fill=fgorange!15] at (\i*0.8,0) {\i};
}
\node[font=\small, anchor=west] at (5.4,1.6) {\texttt{a}: kształt $[2,3]$, kroki $(3,1)$};
\foreach \r in {0,1} \foreach \c in {0,1,2} {
  \pgfmathtruncatemacro{\v}{3*\r+\c}
  \node[cell, fill=fgblue!12] at (5.8+\c*0.8, 0.6-\r*0.8) {\v};
}
\node[font=\small, anchor=west] at (9.4,1.6) {\texttt{a.t()}: kształt $[3,2]$, kroki $(1,3)$};
\foreach \r in {0,1,2} \foreach \c in {0,1} {
  \pgfmathtruncatemacro{\v}{\r+3*\c}
  \node[cell, fill=fggreen!14] at (9.8+\c*0.8, 0.6-\r*0.8) {\v};
}
\end{tikzpicture}
\caption{Element $(i,j)$ widoku leży pod adresem
  $\textit{offset} + i\cdot\textit{stride}_0 + j\cdot\textit{stride}_1$, więc transpozycja
  zamienia tylko kroki, bez kopiowania danych.}
\label{fig:strides}
\end{figure}

\subsection{Jądra CUDA (Faza 2)}

Wszystkie operacje matematyczne mają własne jądra (tab.~\ref{tab:kernels}). Jądra używają
pętli \emph{grid-stride} lub jawnych sprawdzeń granic, 64-bitowych indeksów i~wskaźników
\texttt{\_\_restrict\_\_}; wejścia nieciągłe obsługiwane są przez kroki.

\begin{table}[H]
\centering
\caption{Zaimplementowane operacje i~jądra CUDA.}
\label{tab:kernels}
\begin{tabularx}{\textwidth}{@{}lXl@{}}
\toprule
\textbf{Grupa} & \textbf{Operacje} & \textbf{Plik} \\
\midrule
Jednoargumentowe & \texttt{neg}, \texttt{relu}, \texttt{exp}, \texttt{log}, \texttt{sigmoid}
  (stabilna numerycznie), \texttt{reciprocal} & \texttt{elementwise.cu} \\
Dwuargumentowe & \texttt{add}, \texttt{sub}, \texttt{mul}, \texttt{div}, \texttt{maximum},
  \texttt{lt}, \texttt{eq}; broadcasting, operatory ze skalarami & \texttt{elementwise.cu} \\
W miejscu & \texttt{add\_}, \texttt{sub\_}, \texttt{mul\_}, \texttt{div\_}, \texttt{fill\_},
  kopia z~krokami & \texttt{elementwise.cu} \\
Redukcje & \texttt{sum}, \texttt{prod}, \texttt{max}, \texttt{mean} po wymiarze
  i~po wszystkich elementach; \texttt{unbroadcast} & \texttt{reduce.cu} \\
Softmax & \texttt{log\_softmax} (blok na wiersz, redukcja max i~sumy) & \texttt{softmax.cu} \\
Mnożenie macierzy & \texttt{matmul}: kafelek bloku $64\times64$, krok $K$ równy 8,
  kafelek rejestrowy $8\times8$ na wątek, transpozycja $A$ w~pamięci współdzielonej,
  odczyty \texttt{float4} & \texttt{matmul.cu} \\
\bottomrule
\end{tabularx}
\end{table}

Jądro \texttt{matmul} przepisano w~PR \#12 z~prostego kafelkowania $16\times16$ (jeden
element wyniku na wątek) na wersję z~kafelkowaniem rejestrowym (rys.~\ref{fig:mmtile}).
Każdy z~256 wątków bloku liczy 64 elementy wyniku, co zwiększa liczbę operacji na jeden
odczyt z~pamięci współdzielonej. Jądro wymaga wierszy wyrównanych do 16 bajtów, więc
macierze o~innym układzie są najpierw kopiowane do bufora z~dopełnieniem.

\begin{figure}[H]
\centering
\begin{tikzpicture}[x=0.75cm, y=0.75cm]
% A tile
\draw[fill=fgblue!12] (0,0) rectangle (1,4);
\node[font=\small] at (0.5,4.4) {$A$: $64\times8$};
\draw[fill=fgblue!45] (0,2.5) rectangle (1,3);
% B tile
\draw[fill=fggreen!12] (2,5) rectangle (6,6);
\node[font=\small, anchor=west] at (6.2,5.5) {$B$: $8\times64$};
\draw[fill=fggreen!45] (3.5,5) rectangle (4,6);
% C tile
\draw[fill=fgorange!10] (2,0) rectangle (6,4);
\foreach \i in {1,...,7} {
  \draw[gray!40] (2+\i*0.5,0) -- (2+\i*0.5,4);
  \draw[gray!40] (2,\i*0.5) -- (6,\i*0.5);
}
\draw[fill=fgorange!60] (3.5,2.5) rectangle (4,3);
\node[font=\small, anchor=west] at (6.2,2) {$C$: kafelek bloku $64\times64$};
\node[font=\small, anchor=west, align=left] at (6.2,1.0)
  {pole = kafelek $8\times8$ jednego wątku\\(64 akumulatory w~rejestrach)};
\draw[darr, fgblue] (1,2.75) -- (3.5,2.75);
\draw[darr, fggreen] (3.75,5) -- (3.75,3);
\node[font=\footnotesize, align=center] at (0.5,-0.6) {pamięć\\współdzielona};
\end{tikzpicture}
\caption{Podział pracy w~jądrze \texttt{matmul}: blok wczytuje paski $A$ i~$B$ do pamięci
  współdzielonej, a~każdy wątek liczy kafelek $8\times8$ wyniku.}
\label{fig:mmtile}
\end{figure}

\subsection{Alokator buforujący (Faza 3)}

\texttt{cudaFree} synchronizuje całe urządzenie, a~\texttt{cudaMalloc} jest kosztowne.
Alokator przechowuje zwolnione bloki w~puli posortowanej według rozmiaru
(\texttt{std::multimap}) i~wybiera najmniejszy pasujący blok (\emph{best fit}); nadmiar
jest ograniczony do 16~MiB i~do rozmiaru żądania. Przy braku pamięci zwalnia pulę i~ponawia
próbę. \texttt{FLASHGRAD\_NO\_CACHING=1} wyłącza buforowanie, co jest potrzebne, żeby
\texttt{compute-sanitizer} wykrywał dostęp poza granice bloku.

\begin{figure}[H]
\centering
\begin{tikzpicture}[node distance=6mm and 9mm]
\node[box, fill=fgblue!12] (req) {\texttt{allocate(n)}: zaokrąglij do 512 B};
\node[diamond, draw, aspect=2.4, fill=fgorange!15, below=of req, font=\small,
      inner sep=1pt, align=center] (hit) {blok w~puli\\(best fit)?};
\node[box, fill=fggreen!15, right=16mm of hit] (take) {weź blok z~puli};
\node[box, below=of hit] (mal) {\texttt{cudaMalloc}};
\node[diamond, draw, aspect=2.4, fill=fgorange!15, below=of mal, font=\small,
      inner sep=1pt] (oom) {OOM?};
\node[box, fill=fgred!12, left=14mm of oom] (empty) {\texttt{empty\_cache()}\\i~ponów};
\node[box, fill=fggreen!15, below=of oom] (ret) {zwróć wskaźnik, statystyki};
\node[box, fill=fggray!12, right=16mm of ret] (free) {\texttt{free(p)}: blok wraca do puli};
\draw[arr] (req) -- (hit);
\draw[arr] (hit) -- node[above, font=\small]{tak} (take);
\draw[arr] (hit) -- node[right, font=\small]{nie} (mal);
\draw[arr] (mal) -- (oom);
\draw[arr] (oom) -- node[above, font=\small]{tak} (empty);
\draw[arr] (empty) |- (mal);
\draw[arr] (oom) -- node[right, font=\small]{nie} (ret);
\draw[arr] (take) |- (ret);
\end{tikzpicture}
\caption{Przydział pamięci w~alokatorze buforującym.}
\label{fig:alloc}
\end{figure}

\subsection{Automatyczne różniczkowanie (Faza 4)}

Różniczkowalna operacja liczy wynik, a~następnie -- jeśli któreś wejście wymaga
gradientu -- rejestruje węzeł grafu z~funkcją wsteczną zbudowaną z~innych operacji:

\begin{lstlisting}[style=cpp]
Tensor out = binary(BinaryOp::Mul, "mul", a, b);  // jądro CUDA
if (needs_grad({a, b})) {
  record(out, "mul", {a, b},
         [a = a.detach(), b = b.detach(), need_a = a.requires_grad(),
          need_b = b.requires_grad()](const Tensor& g) {
           return Grads{need_a ? g * b : Tensor(), need_b ? g * a : Tensor()};
         });
}
return out;
\end{lstlisting}

Silnik \texttt{run\_backward} zlicza krawędzie wchodzące do każdego węzła osiągalnego
z~korzenia i~przetwarza węzły w~porządku topologicznym -- węzeł uruchamia się po zebraniu
gradientów od wszystkich konsumentów. Gradienty po broadcastingu są sumowane do kształtu
wejścia (\texttt{unbroadcast}). Funkcje wsteczne działają bez rejestrowania grafu
(\texttt{NoGradGuard}), a~operacje w~miejscu na tensorach wymagających gradientu zgłaszają
błąd.

\begin{figure}[H]
\centering
\begin{tikzpicture}[
  leaf/.style={box, fill=fgblue!12, minimum width=1.2cm},
  op/.style={box, fill=fgorange!15, rounded corners=6pt, minimum width=1.9cm},
  node distance=5mm and 6mm]
\node[leaf] (x) {\texttt{x}};
\node[leaf, right=of x] (w1) {\texttt{W1}};
\node[leaf, right=10mm of w1] (b1) {\texttt{b1}};
\coordinate (xw) at ($(x)!0.5!(w1)$);
\node[op, below=of xw] (mm1) {\texttt{matmul}};
\node[op, below=of mm1] (add1) {\texttt{add}};
\node[op, below=of add1] (relu) {\texttt{relu}};
\node[leaf, right=10mm of relu] (w2) {\texttt{W2}, \texttt{b2}};
\node[op, below=of relu] (mm2) {\texttt{matmul}, \texttt{add}};
\node[op, below=of mm2] (ls) {\texttt{log\_softmax}};
\node[op, below=of ls] (loss) {\texttt{cross\_entropy}};
\node[box, fill=fgred!15, below=of loss] (L) {\texttt{loss}};
\draw[arr] (x) -- (mm1); \draw[arr] (w1) -- (mm1);
\draw[arr] (mm1) -- (add1); \draw[arr] (b1) |- (add1);
\draw[arr] (add1) -- (relu); \draw[arr] (relu) -- (mm2); \draw[arr] (w2) |- (mm2);
\draw[arr] (mm2) -- (ls); \draw[arr] (ls) -- (loss); \draw[arr] (loss) -- (L);
\coordinate (bk) at ($(L.west)+(-1.6,0)$);
\draw[arr, fgred, dashed] (bk) -- (bk |- mm1)
  node[midway, left, align=center, font=\small, text=fgred]
  {\texttt{backward()}\\porządek\\topologiczny};
\node[leaf, minimum height=5mm, minimum width=7mm, right=30mm of ls] (lg1) {};
\node[right=2mm of lg1, font=\small] {liście (parametry)};
\node[op, minimum height=5mm, minimum width=7mm, below=3mm of lg1] (lg2) {};
\node[right=2mm of lg2, font=\small] {węzły z~funkcją wsteczną};
\end{tikzpicture}
\caption{Graf obliczeń rejestrowany przy kroku w~przód sieci dwuwarstwowej.}
\label{fig:graph}
\end{figure}

\subsection{Moduły \texttt{nn}, optymalizatory i dane (Fazy 5--6)}

\begin{itemize}[leftmargin=*]
\item \textbf{\texttt{nn}} -- klasa bazowa \texttt{Module} (nazwane parametry i~podmoduły,
  tryby \texttt{train()}/\texttt{eval()}), \texttt{Linear} (inicjalizacja jak w~PyTorch),
  \texttt{Sequential}, \texttt{MLP}, aktywacje (\texttt{ReLU}, \texttt{Sigmoid},
  \texttt{Tanh}, \texttt{GELU}, \texttt{SiLU}, \texttt{LeakyReLU}, \texttt{Softmax}),
  \texttt{Dropout} oraz straty \texttt{mse\_loss} i~\texttt{cross\_entropy}.
\item \textbf{\texttt{optim}} -- klasa bazowa \texttt{Optimizer}, \texttt{SGD} (momentum,
  \emph{weight decay}) i~\texttt{Adam}/\texttt{AdamW} z~korekcją obciążenia
  momentów, zgodne z~definicjami PyTorch.
\item \textbf{\texttt{data}} -- \texttt{Dataset}, \texttt{TensorDataset} i~\texttt{DataLoader}
  z~mieszaniem, \texttt{drop\_last} i~wątkami roboczymi (\texttt{std::jthread}). Wątki
  budują paczki wyłącznie w~pamięci hosta; kopia na GPU odbywa się w~wątku, który
  przesuwa iterator, więc cała praca CUDA zostaje w~jednym wątku i~jednym strumieniu.
  Kolejność i~zawartość paczek nie zależy od liczby wątków.
\end{itemize}

%=====================================================================
\section{Pierwszy działający przepływ end-to-end}
%=====================================================================

Kamieniem milowym Fazy~6 jest program \texttt{benchmarks/mnist\_mlp.cpp}, który uczy sieć
$784 \to 256 \to 128 \to 10$ (ReLU, \texttt{Dropout} 0,1) na zbiorze MNIST optymalizatorem
Adam (lr $=10^{-3}$, paczki po 128 przykładów, 5 epok, 4 wątki \texttt{DataLoader}).
Zbiór pobiera CMake podczas budowania. Program kończy się błędem, gdy dokładność testowa
spadnie poniżej 97\%. Przepływ jednego kroku pokazuje rys.~\ref{fig:seq}:

\begin{enumerate}[leftmargin=*]
\item Wątki \texttt{DataLoader} składają paczkę $[128, 1, 28, 28]$ w~pamięci hosta.
\item Główny wątek kopiuje paczkę na GPU (alokacja z~puli alokatora).
\item Krok w~przód przez \texttt{Sequential} i~\texttt{cross\_entropy}; każda operacja
  rejestruje węzeł grafu.
\item \texttt{optimizer.zero\_grad()}, \texttt{loss.backward()} -- silnik autogradu
  liczy gradienty wszystkich 6 tensorów parametrów.
\item \texttt{optimizer.step()} aktualizuje parametry w~trybie bez gradientu.
\item Co 100 kroków \texttt{loss.item()} kopiuje stratę na hosta -- to jedyny punkt
  synchronizacji z~GPU w~pętli uczenia.
\end{enumerate}

\begin{figure}[H]
\centering
\begin{tikzpicture}[
  lane/.style={draw, fill=#1!15, minimum width=2.4cm, minimum height=8mm,
               font=\small\bfseries, rounded corners=2pt, align=center},
  msg/.style={-{Stealth[length=2mm]}, thick},
  lbl/.style={font=\footnotesize, fill=white, inner sep=1pt}]
\node[lane=fgpurple] (w) at (0,0) {Wątki\\\texttt{DataLoader}};
\node[lane=fgblue] (h) at (3.6,0) {Pętla uczenia};
\node[lane=fggreen] (a) at (7.2,0) {Autograd /\\\texttt{optim}};
\node[lane=fgorange] (c) at (10.8,0) {Jądra CUDA};
\node[lane=fggray] (m) at (14.0,0) {Alokator};
\foreach \n in {w,h,a,c,m} \draw[gray, dashed] (\n.south) -- ++(0,-9.4);
\draw[msg] (0,-1.1) -- node[lbl, above]{paczka (CPU)} (3.6,-1.1);
\draw[msg] (3.6,-1.8) -- node[lbl, above]{kopia na GPU} (10.8,-1.8);
\draw[msg] (10.8,-2.3) -- node[lbl, above]{\texttt{allocate}} (14.0,-2.3);
\draw[msg] (3.6,-3.2) -- node[lbl, above]{\texttt{model(x)}, \texttt{cross\_entropy}} (10.8,-3.2);
\draw[msg] (10.8,-3.7) -- node[lbl, above]{\texttt{record()}} (7.2,-3.7);
\draw[msg] (3.6,-4.6) -- node[lbl, above]{\texttt{backward()}} (7.2,-4.6);
\draw[msg] (7.2,-5.2) -- node[lbl, above]{funkcje wsteczne} (10.8,-5.2);
\draw[msg] (3.6,-6.2) -- node[lbl, above]{\texttt{step()}} (7.2,-6.2);
\draw[msg] (7.2,-6.8) -- node[lbl, above]{aktualizacja Adam} (10.8,-6.8);
\draw[msg] (10.8,-7.3) -- node[lbl, above]{\texttt{free} $\to$ pula} (14.0,-7.3);
\draw[msg] (3.6,-8.3) -- node[lbl, above]{\texttt{loss.item()} co 100 kroków} (10.8,-8.3);
\draw[msg, dashed] (10.8,-8.9) -- node[lbl, above]{\texttt{float}} (3.6,-8.9);
\draw[thick, fgred, |-|] (1.6,-1.6) -- (1.6,-7.5)
  node[midway, fill=white, font=\footnotesize, text=fgred, align=center]{2345\\kroków};
\end{tikzpicture}
\caption{Diagram sekwencji kroku uczenia w~\texttt{mnist\_mlp}. Wątki robocze
  przygotowują kolejne paczki równolegle z~obliczeniami na GPU.}
\label{fig:seq}
\end{figure}

\begin{lstlisting}[style=term, caption={Wyjście \texttt{./build/benchmarks/mnist\_mlp} (skrócone).}]
training on 60000 images, 469 batches per epoch
epoch 1, step 100: loss 0.284477
...
epoch 1: test accuracy 0.9512
epoch 2: test accuracy 0.9666
epoch 3: test accuracy 0.9706
epoch 4: test accuracy 0.973
epoch 5: test accuracy 0.975
\end{lstlisting}

Dla porównania uruchomiono \texttt{benchmarks/mnist\_mlp\_torch.py} -- ten sam model,
dane, inicjalizację i~hiperparametry w~PyTorch 2.8 (rys.~\ref{fig:mnist},
tab.~\ref{tab:mnist}). Strumienie losowe są różne, więc wyniki zgadzają się statystycznie,
a nie co do bitu.

\begin{figure}[H]
\centering
\begin{tikzpicture}
\begin{axis}[
  width=8.2cm, height=5.6cm, xlabel={krok}, ylabel={strata (paczka)},
  xmin=0, xmax=2400, ymin=0, ymax=0.4, ymajorgrids, grid style={gray!20},
  axis lines*=left, tick label style={font=\footnotesize}, label style={font=\small},
  legend style={font=\footnotesize}, scaled x ticks=false]
\addplot[thick, fgblue, mark=*, mark size=1.2pt] coordinates {
  (100,0.284477) (200,0.344621) (300,0.229386) (400,0.255039) (500,0.17259) (600,0.072674)
  (700,0.138662) (800,0.102527) (900,0.0736849) (1000,0.116442) (1100,0.0938158)
  (1200,0.162074) (1300,0.109056) (1400,0.0528421) (1500,0.0875813) (1600,0.126071)
  (1700,0.0479068) (1800,0.0364544) (1900,0.0944247) (2000,0.0495429) (2100,0.0337923)
  (2200,0.0386056) (2300,0.0455526)};
\addplot[thick, fggray, dashed, mark=square*, mark size=1.2pt] coordinates {
  (100,0.282795) (200,0.298579) (300,0.172548) (400,0.200614) (500,0.156198)
  (600,0.101636) (700,0.0755943) (800,0.151344) (900,0.0632943) (1000,0.0534815)
  (1100,0.0957797) (1200,0.0569592) (1300,0.0716302) (1400,0.0505606) (1500,0.0363149)
  (1600,0.0780344) (1700,0.060559) (1800,0.0583095) (1900,0.0325318) (2000,0.0962047)
  (2100,0.0260842) (2200,0.0190387) (2300,0.0733986)};
\legend{Flashgrad, PyTorch}
\end{axis}
\end{tikzpicture}\hfill
\begin{tikzpicture}
\begin{axis}[
  width=7.4cm, height=5.6cm, xlabel={epoka}, ylabel={dokładność testowa [\%]},
  xmin=0.8, xmax=5.2, ymin=94.5, ymax=98, xtick={1,...,5}, ymajorgrids,
  grid style={gray!20}, axis lines*=left, tick label style={font=\footnotesize},
  label style={font=\small}, legend style={font=\footnotesize, at={(0.97,0.03)},
  anchor=south east}]
\addplot[thick, fgblue, mark=*] coordinates
  {(1,95.12) (2,96.66) (3,97.06) (4,97.30) (5,97.50)};
\addplot[thick, fggray, dashed, mark=square*] coordinates
  {(1,95.25) (2,96.78) (3,97.27) (4,97.69) (5,97.66)};
\addplot[fgred, densely dotted, thick, domain=0.8:5.2] {97};
\legend{Flashgrad, PyTorch, próg 97\%}
\end{axis}
\end{tikzpicture}
\caption{Uczenie MLP na MNIST: strata co 100 kroków (lewy) i~dokładność testowa po
  każdej epoce (prawy).}
\label{fig:mnist}
\end{figure}

\begin{table}[H]
\centering
\caption{MNIST MLP: Flashgrad i~PyTorch z~tą samą konfiguracją.}
\label{tab:mnist}
\begin{tabular}{@{}lrr@{}}
\toprule
 & \textbf{Flashgrad} & \textbf{PyTorch 2.8} \\
\midrule
Dokładność testowa po 5 epokach & 97,50\% & 97,66\% \\
Czas całego programu (wczytanie danych + uczenie + ewaluacja) & \SI{4.2}{s} & \SI{21.0}{s} \\
\bottomrule
\end{tabular}
\end{table}

Różnica czasu wynika głównie z~narzutu Pythona: sam start PyTorch z~inicjalizacją CUDA
trwa ok.\ \SI{1.5}{s}, a~\texttt{DataLoader} PyTorch składa paczki w~Pythonie, przykład po
przykładzie. To porównanie całego programu, a~nie samych jąder -- szczegółowe
profilowanie (Nsight Systems) jest zaplanowane w~Fazie~9.

%=====================================================================
\section{Raport z testów}
%=====================================================================

\subsection{Metodyka}

\begin{itemize}[leftmargin=*]
\item \textbf{Testy z~referencją.} Każda operacja jest porównywana z~naiwną referencją
  w~kodzie testów (\texttt{tests/reference.h}: pętle po \texttt{std::vector},
  akumulatory \texttt{double}) na losowych danych, w~tym: puste tensory, rozmiar 1,
  rozmiary niebędące wielokrotnością bloku (257, 1\,000\,003), wejścia nieciągłe
  i~broadcasting.
\item \textbf{Tolerancje.} Porównanie dokładne tylko tam, gdzie jest gwarantowane
  (pojedyncze operacje fp32, kopie); w~pozostałych przypadkach tolerancja względna
  i~bezwzględna.
\item \textbf{Sprawdzanie gradientów.} Gradienty z~\texttt{backward()} (fp32, GPU)
  porównywane są z~różnicami centralnymi referencji liczonymi w~\texttt{double}
  (\texttt{tests/gradcheck.h}).
\item \textbf{Dwie konfiguracje alokatora.} CTest uruchamia zestaw z~buforowaniem
  i~z~\texttt{FLASHGRAD\_NO\_CACHING=1}.
\item \textbf{Sanitizery.} Kompilacja Debug z~synchronizacją po każdym jądrze uruchamiana
  pod \texttt{compute-sanitizer}.
\end{itemize}

\subsection{Wyniki testów jednostkowych}

\begin{lstlisting}[style=term, caption={Wynik \texttt{ctest --test-dir build} (Release, \texttt{main}, 18.11.2026).}]
974/974 Test #974: no_caching.Workers/DataLoaderWorkers.RejectsCudaExamples/8 ... Passed 0.36 sec

100% tests passed, 0 tests failed out of 974

Total Test time (real) = 132.29 sec

The following tests did not run:
    617 - no_caching.CachingAllocator.ReusesAFreedBlock (Skipped)
    618 - no_caching.CachingAllocator.ReusesTheMostRecentlyFreedBlockOfASize (Skipped)
    619 - no_caching.CachingAllocator.PicksTheSmallestBlockThatFits (Skipped)
    620 - no_caching.CachingAllocator.SmallRequestDoesNotTakeALargeBlock (Skipped)
    624 - no_caching.CachingAllocator.RetriesAfterEmptyingTheCacheWhenOutOfMemory (Skipped)

$ ./build/tests/flashgrad_tests --gtest_brief=1
[==========] 487 tests from 64 test suites ran. (16910 ms total)
[  PASSED  ] 487 tests.
\end{lstlisting}

Pięć pominiętych testów sprawdza ponowne użycie bloków z~puli, więc bez buforowania
celowo wywołują \texttt{GTEST\_SKIP()}.

\begin{figure}[H]
\centering
\begin{tikzpicture}
\begin{axis}[
  xbar, width=12cm, height=10.5cm, bar width=6pt, bar shift=0pt,
  xmin=0, xmax=100, xlabel={liczba testów},
  symbolic y coords={
    adam,sgd,
    dataset,dataloader,
    activation,dropout,loss,linear,module,sequential,functional,
    gradcheck,grad\_mode,in\_place,grad\_softmax,grad\_matmul,grad\_views,grad\_reduce,
    engine,tensor\_autograd,grad\_elementwise,
    error,dtype\_device,shape,views,tensor,
    runtime,ops\_matmul,smoke,ops\_fill,ops\_softmax,caching\_allocator,tensor\_cuda,
    ops\_reduce,ops\_elementwise},
  ytick=data, y tick label style={font=\scriptsize\ttfamily},
  nodes near coords, nodes near coords style={font=\scriptsize},
  axis lines*=left, xmajorgrids, grid style={gray!20},
  legend style={at={(0.97,0.03)}, anchor=south east, font=\footnotesize}]
\addplot[fill=fgblue!60, draw=fgblue!80!black] coordinates {
  (4,runtime) (7,ops\_matmul) (8,smoke) (9,ops\_fill) (9,ops\_softmax)
  (10,caching\_allocator) (21,tensor\_cuda) (29,ops\_reduce) (34,ops\_elementwise)};
\addplot[fill=fgorange!60, draw=fgorange!80!black] coordinates {
  (5,error) (6,dtype\_device) (7,shape) (17,views) (28,tensor)};
\addplot[fill=fggreen!60, draw=fggreen!80!black] coordinates {
  (3,gradcheck) (4,grad\_mode) (4,in\_place) (4,grad\_softmax) (5,grad\_matmul)
  (9,grad\_views) (10,grad\_reduce) (11,engine) (11,tensor\_autograd) (36,grad\_elementwise)};
\addplot[fill=fgpurple!60, draw=fgpurple!80!black] coordinates {
  (9,activation) (7,dropout) (9,loss) (11,linear) (12,module) (13,sequential)
  (17,functional)};
\addplot[fill=fgred!50, draw=fgred!80!black] coordinates {(8,dataset) (92,dataloader)};
\addplot[fill=fggray!50, draw=fggray] coordinates {(8,adam) (10,sgd)};
\legend{\texttt{cuda} (131), \texttt{core} (63), \texttt{autograd} (97), \texttt{nn} (78),
  \texttt{data} (100), \texttt{optim} (18)}
\end{axis}
\end{tikzpicture}
\caption{Liczba testów w~plikach \texttt{tests/<moduł>/test\_<nazwa>.cpp} (łącznie 487).
  Testy \texttt{DataLoader} są sparametryzowane liczbą wątków (0, 1, 2, 4, 8).}
\label{fig:tests}
\end{figure}

\subsection{Pokrycie kodu}

Pokrycie zmierzono osobną kompilacją Debug z~\texttt{--coverage} (dla plików \texttt{.cu}
przez \texttt{-Xcompiler --coverage}) i~narzędziem \texttt{gcov}. \texttt{gcov} nie mierzy
kodu wykonywanego na GPU -- poprawność jąder weryfikują testy z~referencją
i~\texttt{compute-sanitizer}.

\begin{table}[H]
\centering
\caption{Pokrycie linii kodu hosta (gcov, \texttt{main}).}
\label{tab:cov}
\begin{tabular}{@{}lrrr@{}}
\toprule
\textbf{Moduł} & \textbf{Pokryte linie} & \textbf{Linie wykonywalne} & \textbf{Pokrycie} \\
\midrule
\texttt{src/autograd} & 125 & 125 & 100,0\% \\
\texttt{src/core}     & 763 & 778 & 98,1\% \\
\texttt{src/cpu}      & 33  & 33  & 100,0\% \\
\texttt{src/cuda} (część hosta) & 353 & 400 & 88,2\% \\
\texttt{src/data}     & 149 & 151 & 98,7\% \\
\texttt{src/nn}       & 140 & 144 & 97,2\% \\
\texttt{src/optim}    & 60  & 60  & 100,0\% \\
\midrule
\textbf{Razem \texttt{src/}} & \textbf{1623} & \textbf{1691} & \textbf{96,0\%} \\
\bottomrule
\end{tabular}
\end{table}

Niepokryte linie w~\texttt{src/cuda} to głównie obsługa błędów niemożliwych do wywołania
w~teście (np.\ przekroczenie limitu siatki, wymagające tensora większego niż pamięć GPU)
i~ścieżki dla wielu urządzeń.

\subsection{Weryfikacja pamięci i synchronizacji}

\begin{table}[H]
\centering
\caption{Wyniki \texttt{compute-sanitizer} (Debug, \texttt{FLASHGRAD\_NO\_CACHING=1}).}
\label{tab:san}
\begin{tabularx}{\textwidth}{@{}llX@{}}
\toprule
\textbf{Narzędzie} & \textbf{Wynik} & \textbf{Zakres i~uwagi} \\
\midrule
\texttt{memcheck}  & 0 błędów dostępu & Cały zestaw (482 zaliczone, 5 pominiętych). Zgłoszone
  4 błędy API to celowe \texttt{cudaErrorMemoryAllocation} z~testów
  \texttt{OutOfMemory\ldots}, sprawdzających, że OOM zgłasza wyjątek i~zostawia CUDA
  w~stanie używalnym. \\
\texttt{racecheck} & 0 zagrożeń & 31 testów jąder z~pamięcią współdzieloną
  (\texttt{matmul}, redukcje, \texttt{log\_softmax}). \\
\texttt{synccheck} & 0 błędów & Te same 31 testów: poprawne użycie \texttt{\_\_syncthreads()}. \\
\texttt{initcheck} & 0 błędów & Te same 31 testów: brak odczytów niezainicjalizowanej pamięci. \\
\bottomrule
\end{tabularx}
\end{table}

%=====================================================================
\section{Wydajność jąder}
%=====================================================================

Pomiary wykonano na RTX 3090 w~kompilacji Release. Czas mierzono zdarzeniami CUDA
(\texttt{cuda::EventTimer}), uśredniając 10--50 powtórzeń po rozgrzewce. cuBLAS
(\texttt{cublasSgemm}) jest wyłącznie punktem odniesienia. Wartości dla poprzedniego jądra
\texttt{matmul} zmierzono tym samym programem na commicie \texttt{56d0d79}, przed PR \#12.

\begin{figure}[H]
\centering
\begin{tikzpicture}
\begin{axis}[
  width=13cm, height=6.4cm, xmode=log, log basis x=2,
  xtick={128,256,512,1024,2048,4096}, xticklabels={128,256,512,1024,2048,4096},
  xlabel={$n$ (macierze $n\times n$)}, ylabel={GFLOP/s},
  ymin=0, ymax=23000, ymajorgrids, grid style={gray!20}, axis lines*=left,
  legend pos=north west, legend style={font=\footnotesize},
  tick label style={font=\footnotesize},
  scaled y ticks=false, yticklabel style={/pgf/number format/1000 sep={\,}}]
\addplot[thick, fgblue, mark=*] coordinates
  {(128,282.1) (256,1214.5) (512,5025.8) (1024,10453.4) (2048,14506.9) (4096,15716.4)};
\addplot[thick, fgorange, mark=triangle*] coordinates
  {(128,723.7) (256,1713.8) (512,2264.0) (1024,2344.2) (2048,2398.2) (4096,2488.1)};
\addplot[thick, fggray, dashed, mark=square*] coordinates
  {(128,706.2) (256,4025.6) (512,10922.7) (1024,18648.0) (2048,22788.9) (4096,21508.9)};
\legend{Flashgrad \texttt{matmul} (register tiling; \texttt{main}),
  poprzednie jądro (tiling $16\times16$),
  cuBLAS \texttt{sgemm}}
\end{axis}
\end{tikzpicture}
\caption{Przepustowość mnożenia macierzy fp32.}
\label{fig:matmul}
\end{figure}

Nowe jądro osiąga 15,7~TFLOP/s dla $n=4096$, czyli 73\% wydajności cuBLAS i~6,3 razy
więcej niż poprzednie. Dla małych macierzy ($n \le 256$) jest jednak wolniejsze od
poprzedniego, bo przygotowanie buforów z~dopełnieniem i~dodatkowe kopie kosztują więcej
niż samo mnożenie. Ma to znaczenie dla małych sieci (tab.~\ref{tab:cache}) i~jest
zapisane jako zadanie do poprawy.

\begin{figure}[H]
\centering
\begin{tikzpicture}
\begin{axis}[
  width=13cm, height=6cm, xmode=log, log basis x=2,
  xtick={65536,1048576,4194304,16777216,67108864},
  xticklabels={$2^{16}$,$2^{20}$,$2^{22}$,$2^{24}$,$2^{26}$},
  xlabel={liczba elementów}, ylabel={przepustowość [GB/s]},
  ymin=0, ymax=1000, ymajorgrids, grid style={gray!20}, axis lines*=left,
  legend style={font=\footnotesize, at={(0.98,0.02)}, anchor=south east},
  tick label style={font=\footnotesize}]
\addplot[thick, fgblue, mark=*] coordinates
  {(65536,243.8) (1048576,742.5) (4194304,805.8) (16777216,712.0) (67108864,804.8)};
\addplot[thick, fgblue, dashed, mark=o] coordinates
  {(65536,79.6) (1048576,79.3) (4194304,104.6) (16777216,135.7) (67108864,253.5)};
\addplot[thick, fggreen, mark=square*] coordinates
  {(65536,49.2) (1048576,572.9) (4194304,670.1) (16777216,799.7) (67108864,738.4)};
\addplot[thick, fgred, densely dotted, domain=65536:67108864] {936};
\legend{\texttt{a + b} z~buforowaniem, \texttt{a + b} bez buforowania,
  \texttt{sum(a)} z~buforowaniem, szczytowa przepustowość (936 GB/s)}
\end{axis}
\end{tikzpicture}
\caption{Przepustowość pamięci jąder elementwise i~redukcji (fp32).}
\label{fig:bw}
\end{figure}

Dodawanie osiąga do 806~GB/s (86\% szczytowej przepustowości GDDR6X), a~\texttt{sum} do
800~GB/s (85\%), więc jądra ograniczone pamięcią są blisko limitu sprzętowego. Bez
alokatora buforującego przepustowość dodawania spada 3--9-krotnie, bo każdy wynik wymaga
\texttt{cudaMalloc}, a~każde zwolnienie -- synchronizującego \texttt{cudaFree}.

\begin{table}[H]
\centering
\caption{Wpływ alokatora buforującego.}
\label{tab:cache}
\begin{tabular}{@{}lrrr@{}}
\toprule
\textbf{Pomiar} & \textbf{z buforowaniem} & \textbf{bez buforowania} & \textbf{przyspieszenie} \\
\midrule
\texttt{a + b}, $2^{20}$ el.     & \SI{0.017}{ms} & \SI{0.159}{ms} & 9,4$\times$ \\
\texttt{sum(a)}, $2^{20}$ el.    & \SI{0.007}{ms} & \SI{0.016}{ms} & 2,2$\times$ \\
\texttt{matmul} $512\times512$   & \SI{0.053}{ms} & \SI{0.094}{ms} & 1,8$\times$ \\
krok uczenia sieci $2\to32\to2$ (w~przód + wstecz) & \SI{0.179}{ms} & \SI{0.343}{ms} & 1,9$\times$ \\
\bottomrule
\end{tabular}
\end{table}

\texttt{DataLoader} składa cały zbiór treningowy MNIST (60\,000 obrazów) w~\SI{0.274}{s}
w~jednym wątku i~w~\SI{0.113}{s} z~4 wątkami roboczymi (2,4$\times$ szybciej;
\texttt{examples/dataloader\_basics}).

%=====================================================================
\section{Zgodność z harmonogramem i ścieżka krytyczna}
%=====================================================================

\begin{figure}[H]
\centering
\begin{tikzpicture}[x=0.2cm, y=0.6cm]
% dzień 0 = 13.10.2026; 110 = 31.01.2027
\foreach \d/\l in {0/13.10, 19/1.11, 49/1.12, 80/1.01, 110/31.01} {
  \draw[gray!40] (\d,0.5) -- (\d,-11.4);
  \node[font=\scriptsize, above] at (\d,0.5) {\l};
}
\draw[fgred, thick] (36,0.5) -- (36,-11.4) node[below, font=\scriptsize, text=fgred] {dziś (18.11)};
\newcommand{\fgtask}[5]{% wiersz, początek, koniec, kolor, etykieta
  \node[font=\scriptsize, anchor=east] at (-1,-#1) {#5};
  \fill[#4, rounded corners=1pt] (#2,-#1-0.32) rectangle (#3,-#1+0.32);}
\fgtask{0}{0}{1}{fggreen}{F0 infrastruktura}
\fgtask{1}{2}{4}{fggreen}{F1 rdzeń tensorów}
\fgtask{2}{4}{8}{fggreen}{F2 jądra CUDA}
\fgtask{3}{8}{18}{fggreen}{F2 \texttt{matmul} (PR \#12)}
\fgtask{4}{6}{8}{fggreen}{F3 alokator}
\fgtask{5}{7}{17}{fggreen}{F4 autograd}
\fgtask{6}{29}{34}{fggreen}{F5 moduły \texttt{nn}}
\fgtask{7}{31}{34}{fggreen}{F6 optim, dane, MNIST}
\fgtask{8}{37}{56}{fgblue!50}{F7 Conv2d, CIFAR-10}
\fgtask{9}{56}{96}{fgblue!50}{F8 FlashAttention, GPT}
\fgtask{10}{90}{110}{fgblue!50}{F9 benchmarki}
\end{tikzpicture}
\caption{Harmonogram faz (zielone -- zakończone, według historii git; niebieskie --
  planowane).}
\label{fig:gantt}
\end{figure}

\begin{table}[H]
\centering
\caption{Stan zadań Etapu 3.}
\label{tab:schedule}
\begin{tabularx}{\textwidth}{@{}lcX@{}}
\toprule
\textbf{Zadanie} & \textbf{Status} & \textbf{Uwagi i~ewentualne opóźnienia} \\
\midrule
F0--F1: infrastruktura, tensory & \textbf{100\%} & PR \#1--\#6, zgodnie z~planem \\
F2: jądra CUDA & \textbf{100\%} & PR \#8, \#9, \#11, \#12; \texttt{matmul} 73\% cuBLAS \\
F3: alokator buforujący & \textbf{100\%} & PR \#10; do 9,4$\times$ szybsze małe operacje \\
F4: autograd & \textbf{100\%} & PR \#13, \#14; gradienty sprawdzone różnicami skończonymi \\
F5: moduły \texttt{nn} & \textbf{100\%} & PR \#16; \texttt{mlp\_basics}: 97,9\% dokładności \\
F6: optymalizatory, \texttt{DataLoader}, MNIST & \textbf{100\%} & PR \#15, \#17; MNIST
  97,5\%. Do zrobienia: odhaczenie fazy w~\texttt{README.md} \\
Narzut \texttt{matmul} dla małych macierzy & \textbf{otwarte} & Regresja dla
  $n \le 256$; nie blokuje F7 \\
\bottomrule
\end{tabularx}
\end{table}

\paragraph{Ścieżka krytyczna.} Fazy 0--6 zakończono przed terminem, co daje zapas na
najtrudniejszą Fazę~8 (FlashAttention, mieszana precyzja, mały GPT), zależną od jąder
złożonych i~\texttt{Conv2d} z~Fazy~7. Główne ryzyka to wprowadzenie fp16/bf16 (rozszerzenie
\texttt{DType} i~wszystkich jąder) oraz narzut małych operacji, który w~modelu GPT
z~wieloma małymi macierzami może dominować. Projekt jest na to przygotowany: tensory od
początku przechowują typ danych, a~API jąder przyjmuje strumień CUDA.

\paragraph{Następne kroki (Etap 4).} \texttt{Conv2d} i~pooling z~własnymi jądrami, jądra
złożone (np.\ \texttt{linear}+aktywacja), uczenie CNN na CIFAR-10, usunięcie narzutu
\texttt{matmul} dla małych macierzy.

\vfill
\noindent
\begin{tabularx}{\textwidth}{@{}Xr@{}}
\textbf{Podpis promotora (akceptacja Etapu 3):} \dotfill & \textbf{Data:} 22 listopada 2026 \\
\end{tabularx}

\end{document}
